\paragraph{Author oracle-machine conventions (Sections 2.2--2.3).}
We assume without loss of generality that a polynomial-time bounded randomized Turing Machine $M$ comes equipped with a poly-time computable function $r^M:\{0,1\}^*\to\mathbb N$ where $r_n^M$ is a polynomial upper bound on the number of random bits generated by the machine on input $x$. On input $x$ the TM starts by generating an auxiliary random string $\rho$ of length $r_x^M$ which is written on the random tape, and the program operates deterministicially on the pair $x,\rho$.
Normally, an oracle is a function $O$ from a query set $Q$ to an answer set $A$. We also need a more general notion of a context-sensitive oracle in which the response depends on the query $q$ and the context (machine $M$ and input $x$ that is calling the oracle). We write $O_x^M(q)$ for the output of oracle $O$ to query $q$.
Unless otherwise specified all oracles in this paper have query set $Q=\{0,1\}^*$ and answer set $A=\mathbb N$.
An OTM $M$ instantiated by oracle $O$, denoted $M^O$ is the program that uses oracle $O$ to answer any queries. If $O$ is computable, then $M^O$ is computably instantiated, and can be identified with an oracle-free Turing machine, even in the case that $O$ is a context-sensitive oracle.
The randomized OTMs we consider are decision machines which means that they always halt and output either 1 (identified with accept) or 0 (identified with reject). The language $L(M^O)$ recognized by an (possibly randomized) instantiated decision OTM $M^O$ is the set of $x$ for which the probability that $M(x)$ accepts is greater than $1/2$.
Let $\epsilon:\mathbb N\to[0,1]$. We say that $L$ is accepted by $M^O$ with advantage $\epsilon$ if every accepted string $x$ is accepted with probability at least $(1+\epsilon(|x|))/2$ and every rejected string is accepted with probability at most $(1-\epsilon(|x|))/2$. Note that computing with advantage $\epsilon=1$ means that $M^O(x)$ always outputs $L(x)$.
An algorithm that accepts $L$ with advantage $\epsilon$ can be converted to an algorithm with advantage $(1-\epsilon')$ by repeating the algorithm $O(\log(1/\epsilon')+(1/\epsilon))^2$ times and taking majority vote.
We restrict attention to OTMs that are nonadaptive. A nonadaptive OTM comes equipped with:
\begin{itemize}
\item A computable function $k^M:\{0,1\}^*\to\mathbb N$. For $x\in\{0,1\}^*$, $k_x^M$ is the number of oracle queries that $M$ makes on input $x$. For convenience (and with no significant loss of generality) we assume that $k_n^M$ is a power of 2.
\item A computable function $q^M$ that takes as input a string $x$ and an integer $i\in\{1,\ldots,k_{|x|}^M\}$, and (if the OTM is randomized) a random string $\rho$ of length $r_{|x|}^M$, and outputs an element of $\{0,1\}^*$. $q_x^M(i,\rho)$ is the $i$th query asked by $M$ on input $x$ when the random string is $\rho$.
\item A computable function $\operatorname{out}^M$ that takes as input $x$ and the sequence of query answers $(a_1,\ldots,a_{k_x^M})$ and (if the OTM is randomized) the random string $\rho$ and outputs either 0 or 1. $\operatorname{out}_x^M(a_1,\ldots,a_{k_x^M},\rho)$ is the output of the algorithm on input $x$, random string $\rho$ when the query answers are $r_1,\ldots,r_{k_x^M}$.
\end{itemize}
In the above definition the nonadaptive OTM $M$ is deterministic if $r_x^M$ is identically 0. Thus the function $q_x^M$ depends only on the query index $i$.
$Q_x^M$ denotes the set of all $q\in\{0,1\}^*$ such that on input $x$, $M$ asks query $q$ with positive probability. $Q_x^M$ is a finite set of size at most $2^{r_x^M}k_x^M$ since there are $2^{r_x^M}$ choices for $\rho$ and each results in $k_x^M$ queries.
We say that $M$ is $\alpha$-honest, for $\alpha:\mathbb N\to\mathbb N$, provided that on any input $x$, all oracle queries have length at least $\alpha(|x|)$. If $\alpha(n)=n^\tau$ for some constant $\tau>0$ we say that $M$ is polynomially honest. Conventionally, a polynomially honest oracle algorithm is referred to simply as honest, but it will be useful for us to distinguish different kinds of honesty.
If $M$ is $\alpha$-honest for some $\alpha$ that tends to $\infty$ we say that $M$ is weakly honest. We say that $M$ is fixed query length if the length of each query of $M$ on input $x$ depends only on $|x|$.
\paragraph{2.3. Approximate Oracles and Robust Oracle Algorithms.}
Recall that in this paper we consider oracles that take as input a string and output an integer. We now introduce two related concepts for an oracle (possibly context-sensitive) $O$:
\begin{itemize}
\item Approximation of $O$ by another oracle $O'$ that is ``suitably close''.
\item An oracle algorithm $M$ instantiated by $O$ that is robust with respect to replacing $O$ by any $O'$ that is close enough to $O$.
\end{itemize}
The closeness of two oracles is measured by a nonnegative valued closeness function $\beta$ with two arguments, a string $q$ and natural number $n$. The value of $\beta$ at these arguments is written $\beta_n(q)$. We say that $O$ and $P$ are $\beta$-close with respect to $M$ provided that for all inputs $x$ to $M$, and queries $q\in Q_x^M$, $|O_x^M(q)-P_x^M(q)|\leq\beta_{|x|}(q)$.
Let $M$ be a randomized oracle algorithm, $O$ be an oracle and $L$ be the language accepted by $M^O$. Let $\beta$ be a closeness function and let $\epsilon:\mathbb N\to[0,1]$. We say that $M$ is $(\beta,\epsilon)$-robust with respect to oracle $O$ provided that for every oracle $O'$ that is $\beta$-close to $O$, $M^{O'}$ accepts $L$ with advantage $\epsilon$. We say that $M$ is $\beta$-robust if it is $(\beta,1)$-robust, which means that for every oracle $O'$ that is $\beta$-close to $O$, $M^{O'}(x)$ always outputs the correct answer $L(x)$.
\paragraph{Proposition 7.}
For every computable function $f:\{0,1\}^*\to\{0,1\}^*$ there is a constant $C_f$ such that if $q=f(z)$ then $K(q)\leq K(z)+C_f$.
In particular, for $f$ being the identity function $\mathsf{ID}$ we have
\paragraph{Proposition 8.}
For all strings $q$, $K(q)\leq|q|+C_{\mathsf{ID}}$.
\paragraph{Proposition 9.}
$\sum_{q\in\{0,1\}^*}2^{-K(q)}/K(q)^2\leq2$.
\begin{proof}
Letting $S_k$ be the set of strings of Kolmogorov complexity exactly $k$, we can rewrite the above sum as:
\[\sum_{k\geq1}\sum_{q\in S_k}\frac{2^{-k}}{k^2}.\]
Since $|S_k|\leq2^k$, this is at most $\sum_{k\geq1}1/k^2<2$.
\end{proof}
The canonical order on $\{0,1\}^*$ is the order in which $x$ precedes $y$ if $|x|<|y|$ or $|x|=|y|$ and $x$ is lexicographically less than $y$.
\paragraph{Proposition 10.}
Let $L$ be a computable language that contains infinitely many strings.
\begin{enumerate}
\item Let $\Lambda:\mathbb N\to\mathbb N$ be any computable unbounded function. Let $x_j$ be the first string (according to the canonical order) satisfying $|x_j|\geq\Lambda(j)$ and $x\in L$. There is an integer $C$ (depending on $L$ and $\Lambda$) such that for all integers $j$, $K(x_j)\leq\log(j)+C$.
\item Let $\kappa:\mathbb N\to\mathbb Z$ be any computable nondecreasing function that tends to $\infty$. Let $L$ be a computable language that contains infinitely many strings. Then $L$ contains a string $x$ for which $K(x)<\kappa(|x|)$.
\end{enumerate}
\begin{proof}
Let $M$ be a Turing Machine that computes $L$ and $N$ be a Turing machine that computes $\Lambda$. We can construct a machine $M'$ that on input $j$ outputs $x_j$. $M$ first evaluates $\Lambda(j)$ and then enumerates strings $x$ of length at least $\Lambda(j)$ in the canonical order and stops when it finds the first string in $L$. $M'$ is a program whose length is a constant (depending on $|N|$ and $|M|$) and so the pair $M',j$ has a description of length $\log(j)+C$.
For the second part, given $\kappa$, let $\Lambda$ be the function where $\Lambda(j)$ is the least integer $m$ such that $2\log(j)<\kappa(m)$. The function $\Lambda$ is computable. Applying the first part, there is a constant $C$ (depending on $L$ and $\Lambda$) such that for all $j$, $x_j$ is a string of length at least $\Lambda(j)$ such that $K(x_j)\leq\log(j)+C$. Choose $j=2^C$. Then $x_j$ has length at least $\Lambda(j)$, and $K(x_j)\leq\log(j)+C\leq2\log(j)<\kappa(\Lambda(j))\leq\kappa(|x_j|)$, as required.
\end{proof}
